\subsection{Outline of the article}
My aims in this paper are both conceptual and empirical.  With this in mind, I set myself three main goals.  First, in section~\ref{sec:establishing}, I attempt to ``establish the phenomenon'' \citep{merton87}, in question by explicitly characterizing the outcome of interest. I show that ``culture-talk'' has been a central theme of various \emph{theoretical} discussions and classic programmatic statements in the culture and stratification literature but that, as noted above, it has not received any sort of sustained empirical attention.  Here I define culture-talk as a species of sociability in the Simmelian sense, argue for the privileged status of aesthetic consumption experiences in this process and advocate for the development of a empirical program of research that focuses on  \emph{the stratification of sociability} in contemporary societies. Second, in section~\ref{sec:importance}, I demonstrate the explanatory potential of the culture-talk concept by considering extant literature on the connection between culture, stratification and social capital with an eye towards showing that culture-talk is a strategic interactional \emph{mechanism} underlying long-standing empirical generalizations in the literature connecting inequalities in social connectivity to social position and culture \citep{erickson08}.  In section~\ref{sec:analysis}, I follow with an empirical analysis aimed at establishing what the socio-demographic predictors of culture-talk in the United States today are.  Given the fact that this is the first paper that attempts to shed empirical light on this phenomenon, I keep the empirical analysis as straightforward as possible---within the limits set by the data at hand---and restrict myself to establishing some key (descriptive) empirical generalizations that I believe are of strategic substantive importance.  I close in section~\ref{sec:discussion} by outlining the implications of my theoretical synthesis and the empirical results for future research in the culture and stratification literature.